% Options for packages loaded elsewhere
\PassOptionsToPackage{unicode}{hyperref}
\PassOptionsToPackage{hyphens}{url}
\PassOptionsToPackage{dvipsnames,svgnames,x11names}{xcolor}
%
\documentclass[
  10pt,
]{article}
\usepackage{amsmath,amssymb}
\usepackage{iftex}
\ifPDFTeX
  \usepackage[T1]{fontenc}
  \usepackage[utf8]{inputenc}
  \usepackage{textcomp} % provide euro and other symbols
\else % if luatex or xetex
  \usepackage{unicode-math} % this also loads fontspec
  \defaultfontfeatures{Scale=MatchLowercase}
  \defaultfontfeatures[\rmfamily]{Ligatures=TeX,Scale=1}
\fi
\usepackage{lmodern}
\ifPDFTeX\else
  % xetex/luatex font selection
  \setmainfont[]{DejaVu Serif}
  \setsansfont[]{DejaVu Sans}
  \setmonofont[]{DejaVu Sans Mono}
  \setmathfont[]{Latin Modern Math}
\fi
% Use upquote if available, for straight quotes in verbatim environments
\IfFileExists{upquote.sty}{\usepackage{upquote}}{}
\IfFileExists{microtype.sty}{% use microtype if available
  \usepackage[]{microtype}
  \UseMicrotypeSet[protrusion]{basicmath} % disable protrusion for tt fonts
}{}
\makeatletter
\@ifundefined{KOMAClassName}{% if non-KOMA class
  \IfFileExists{parskip.sty}{%
    \usepackage{parskip}
  }{% else
    \setlength{\parindent}{0pt}
    \setlength{\parskip}{6pt plus 2pt minus 1pt}}
}{% if KOMA class
  \KOMAoptions{parskip=half}}
\makeatother
\usepackage{xcolor}
\usepackage[margin=2.0cm]{geometry}
\usepackage{longtable,booktabs,array}
\usepackage{calc} % for calculating minipage widths
% Correct order of tables after \paragraph or \subparagraph
\usepackage{etoolbox}
\makeatletter
\patchcmd\longtable{\par}{\if@noskipsec\mbox{}\fi\par}{}{}
\makeatother
% Allow footnotes in longtable head/foot
\IfFileExists{footnotehyper.sty}{\usepackage{footnotehyper}}{\usepackage{footnote}}
\makesavenoteenv{longtable}
\setlength{\emergencystretch}{3em} % prevent overfull lines
\providecommand{\tightlist}{%
  \setlength{\itemsep}{0pt}\setlength{\parskip}{0pt}}
\setcounter{secnumdepth}{-\maxdimen} % remove section numbering
% doc_header.tex -- LaTeX preamble additions used by build_docs.sh for every
% Markdown -> LaTeX -> PDF conversion (supplements, README, submission files).
\usepackage{etoolbox}
% longtable may not be loaded by pandoc when a document has no table
\usepackage{longtable}
% bold vectors: \bm -> \symbf under unicode-math (XeLaTeX), bm package otherwise
\makeatletter
\@ifpackageloaded{unicode-math}{\providecommand{\bm}[1]{\symbf{#1}}}{\usepackage{bm}}
\makeatother
% smaller type in long tables so that wide audit tables fit the text width;
% inside a table a line break is allowed after an underscore of plain text
% (pandoc writes it as \_), so that long identifiers wrap inside their column
% rows are set a little apart, so that symbols with indices do not touch
\AtBeginEnvironment{longtable}{\footnotesize\renewcommand{\arraystretch}{1.45}%
  \let\sptunderscore\_\def\_{\sptunderscore\allowbreak}}
\setlength{\LTcapwidth}{\textwidth}
% allow line breaks after slashes and hyphens inside table cells
\usepackage{xurl}
% code blocks: break over-long lines instead of running into the margin.
% Plain ``` blocks are typeset by pandoc as {verbatim}, highlighted ones
% (```bash) as {Highlighting}; both are re-declared as breakable fancyvrb
% environments.  The continuation line is marked with a small hooked arrow.
\usepackage{fvextra}
\RecustomVerbatimEnvironment{verbatim}{Verbatim}{breaklines,breakanywhere,fontsize=\small}
\makeatletter
\@ifundefined{Highlighting}{}{%
  \RecustomVerbatimEnvironment{Highlighting}{Verbatim}{breaklines,breakanywhere,fontsize=\small,commandchars=\\\{\}}}
\makeatother
% inline code (pandoc writes it as \texttt{...}): never hyphenate -- an inserted
% hyphen could be read as part of a file name -- but allow a line break after
% an underscore, slash, comma, colon or equals sign, so that long paths and
% column lists wrap inside the text width and inside table columns.
% Inside a formula (e.g. the name of a notebook symbol in an equation) the same
% command must first leave math mode, which \text of amsmath does.
\ExplSyntaxOn
\tl_new:N \l__spt_code_tl
\cs_new_protected:Npn \__spt_texttt:n #1
  {
    \tl_set:Nn \l__spt_code_tl {#1}
    \regex_replace_all:nnN { ([\/\,\:\=]) } { \1 \c{allowbreak} } \l__spt_code_tl
    \regex_replace_all:nnN { \c{_} } { \c{_} \c{allowbreak} } \l__spt_code_tl
    \group_begin: \ttfamily \hyphenchar \font = -1 \scan_stop: \l__spt_code_tl \group_end:
  }
\RenewDocumentCommand \texttt { m }
  {
    \mode_if_math:TF { \text { \__spt_texttt:n {#1} } } { \__spt_texttt:n {#1} }
  }
\ExplSyntaxOff
% A figure is written in the Markdown source as an image paragraph followed by
% a caption paragraph ("**Figure S1.** ..."), so that the caption is also
% visible on GitHub.  An image that starts a paragraph is centered.
\usepackage{graphicx}% pandoc loads it only for documents that contain an image
\NewCommandCopy\sptincludegraphics\includegraphics
\RenewDocumentCommand\includegraphics{O{} m}{%
  \ifvmode\begin{center}\sptincludegraphics[#1]{#2}\end{center}%
  \else\sptincludegraphics[#1]{#2}\fi}
% keep figures where they are referenced
\usepackage{float}
\floatplacement{figure}{H}
% running header with the document title
\usepackage{fancyhdr}
\pagestyle{fancy}
\fancyhf{}
\fancyhead[L]{\small PROJECT spacetimes, Revision\_Deflation\_CJT\_and\_Kohn-Sham}
\fancyhead[R]{\small\thepage}
\renewcommand{\headrulewidth}{0.4pt}
% 16x16 matrices: tighter columns so that a full matrix fits the text width
\setlength{\arraycolsep}{3.2pt}
% The documents of this project contain long tagged symbols, file paths and
% identifiers that cannot be hyphenated.  A paragraph is therefore allowed to
% have looser lines before TeX gives up and lets a line stick out of the margin.
\AtBeginDocument{\tolerance=2000 \emergencystretch=5em \hbadness=2000}
\ifLuaTeX
  \usepackage{selnolig}  % disable illegal ligatures
\fi
\IfFileExists{bookmark.sty}{\usepackage{bookmark}}{\usepackage{hyperref}}
\IfFileExists{xurl.sty}{\usepackage{xurl}}{} % add URL line breaks if available
\urlstyle{same}
\hypersetup{
  pdftitle={The task, the rules, what was done, what was not, and the decisions that are the author's},
  pdfauthor={Prepared for Patrick Lee Nash, Ph.D.},
  colorlinks=true,
  linkcolor={blue},
  filecolor={Maroon},
  citecolor={blue},
  urlcolor={blue},
  pdfcreator={LaTeX via pandoc}}

\title{The task, the rules, what was done, what was not, and the
decisions that are the author's}
\usepackage{etoolbox}
\makeatletter
\providecommand{\subtitle}[1]{% add subtitle to \maketitle
  \apptocmd{\@title}{\par {\large #1 \par}}{}{}
}
\makeatother
\subtitle{Provenance document \texttt{P0} of PROJECT spacetimes,
Revision\_Deflation\_CJT\_and\_Kohn-Sham}
\author{Prepared for Patrick Lee Nash, Ph.D.}
\date{November 13, 2026}

\begin{document}
\maketitle

\hypertarget{what-this-document-is}{%
\section{1. What this document is}\label{what-this-document-is}}

This document is the account that the task specification requires of the
work in the folder \texttt{Revision\_Deflation\_CJT\_and\_Kohn-Sham}. It
states

\begin{itemize}
\tightlist
\item
  the task, in the words of this provenance, with the author's
  hypotheses in his formulas (Section 2);
\item
  each of his hard rules and how it is met, with the place where that
  can be verified (Section 3);
\item
  what the folder contains (Section 4) and what was done, in order
  (Section 5);
\item
  the results in brief (Section 6);
\item
  every \textbf{deviation} from the instructions, with its reason
  (Section 7);
\item
  the \textbf{decisions that are the author's} to take or to confirm
  (Section 8), and the findings that concern his own documents (Section
  9);
\item
  what has \textbf{not} been done (Section 10);
\item
  the inputs (Section 11) and the \textbf{log of errors} made during the
  work (Section 12).
\end{itemize}

The three provenance documents of the folder are called \texttt{P0}
(this one), \texttt{P1} (the symbols) and \texttt{P2} (the
calculations), after the beginning of their file names. Every symbol
that occurs here is explained in Section 13. Nothing in this document
has to be looked up elsewhere; where it names another document of the
folder, it does so to say what that document contains, not to define
anything used here.

\hypertarget{the-task}{%
\section{2. The task}\label{the-task}}

\hypertarget{the-hypotheses}{%
\subsection{2.1 The hypotheses}\label{the-hypotheses}}

PROJECT spacetimes studies a sixteen-component Dirac field on two
eight-dimensional spacetimes with the coordinates \(x0,\dots,x7\):
\(x1,x2,x3\) are ordinary space, \(x4\) is the ordinary time, \(x0\) is
a hidden space-like coordinate, and \(x5,x6,x7\) are extra coordinates,
time-like on the \textbf{4+4} spacetime and space-like on the
\textbf{7+1} spacetime.

The author's hypotheses for this sub-project are:

\begin{enumerate}
\def\labelenumi{\arabic{enumi}.}
\tightlist
\item
  \textbf{Deflation.} Each of \(x5\), \(x6\), \(x7\) deflates
  exponentially. Flat spacetime is replaced by a pseudo-Riemannian space
  with a diagonal metric whose components along these coordinates are,
  with three positive rates \(\mathrm{g55},\mathrm{g66},\mathrm{g77}\)
  and positive constants that may differ,
  \[\text{4+4:}\quad g^{(4+4)}_{tt}=-\,\mathrm{const}_t^{\,2}\,e^{-2\,\mathrm{rate}_t\,x4},\qquad
  \text{7+1:}\quad g^{(7+1)}_{tt}=+\,\mathrm{const}_t^{\,2}\,e^{-2\,\mathrm{rate}_t\,x4},\qquad t=5,6,7,\]
  where \(\mathrm{rate}_5=\mathrm{g55}\),
  \(\mathrm{rate}_6=\mathrm{g66}\), \(\mathrm{rate}_7=\mathrm{g77}\).
\item
  \textbf{The Dirac Lagrangians contain a mass term and a
  self-interaction.} The self-interaction is the quartic term of the
  author's document on the 7+1 problem,
  \[\mathcal L_{\rm self}=-\tfrac12\sum_{r=0}^{8}g_r\,I_r,\qquad
  I_r=\sum_{|K|=r}\big(\bar\psi\,\Theta_K\,\psi\big)\big(\bar\psi\,\Theta^K\,\psi\big),\]
  with nine couplings \(g_r\) and the 256 products \(\Theta^K\) of Dirac
  matrices (Section 13). The first revision of this sub-project had
  omitted it.
\item
  \textbf{Two methods.} The effective potential for composite operators
  of Cornwall, Jackiw and Tomboulis (CJT) and the Kohn--Sham ansatz are
  both to be employed: to calculate additional predictions of the two
  models, and as a check on the existing computations. The author's
  solution of the 7+1 problem is to be carried over to the 4+4 problem.
\end{enumerate}

Every calculation of the first revision is to be repeated in full under
these hypotheses, in a new folder, with commented source code,
instructions and provenance for every computed value, and all of the
code is to be executed.

\hypertarget{the-criticisms-that-the-task-states-and-what-was-done-about-each}{%
\subsection{2.2 The criticisms that the task states, and what was done
about
each}\label{the-criticisms-that-the-task-states-and-what-was-done-about-each}}

\begin{longtable}[]{@{}
  >{\raggedright\arraybackslash}p{(\columnwidth - 2\tabcolsep) * \real{0.5000}}
  >{\raggedright\arraybackslash}p{(\columnwidth - 2\tabcolsep) * \real{0.5000}}@{}}
\toprule\noalign{}
\begin{minipage}[b]{\linewidth}\raggedright
criticism
\end{minipage} & \begin{minipage}[b]{\linewidth}\raggedright
what was done
\end{minipage} \\
\midrule\noalign{}
\endhead
\bottomrule\noalign{}
\endlastfoot
The manuscripts are not useful for students: they send the reader to
other manuscripts, documents, provenances or symbol files. & Each
manuscript now contains every symbol, definition and equation that it
uses, with appendices for the symbols, the printed Dirac matrices, the
curvature, and the programs and checks. A script searches every document
for wordings that send the reader elsewhere
(\texttt{check\_self\_contained.py}); none is found. The first revision
had five supplements to which its manuscripts referred; their content is
now inside the manuscripts (Section 7, item V1). \\
The manuscripts do not define the Dirac Lagrangians and Dirac equations
correctly, because the self-interaction is missing. & The Lagrangians
are \(\mathcal L_{\rm kin}+\mathcal L_{\rm mass}+\mathcal L_{\rm self}\)
on both spacetimes, and the field equations carry the cubic term;
manuscripts, Secs. III and IV, and the document ``dirac matrices
deflating'', Parts 3 to 6. The consequences of the two new terms run
through every later section. \\
The same symbol is used for different objects. & A MASTER registry of
229 symbols with one meaning each, the list of the collisions found and
their resolution, and seven machine checks: document \texttt{P1}. All
documents of the folder use it. \\
\end{longtable}

\hypertarget{the-hard-rules-and-how-each-is-met}{%
\section{3. The hard rules, and how each is
met}\label{the-hard-rules-and-how-each-is-met}}

\begin{longtable}[]{@{}
  >{\raggedright\arraybackslash}p{(\columnwidth - 4\tabcolsep) * \real{0.3333}}
  >{\raggedright\arraybackslash}p{(\columnwidth - 4\tabcolsep) * \real{0.3333}}
  >{\raggedright\arraybackslash}p{(\columnwidth - 4\tabcolsep) * \real{0.3333}}@{}}
\toprule\noalign{}
\begin{minipage}[b]{\linewidth}\raggedright
rule
\end{minipage} & \begin{minipage}[b]{\linewidth}\raggedright
how it is met
\end{minipage} & \begin{minipage}[b]{\linewidth}\raggedright
where to verify
\end{minipage} \\
\midrule\noalign{}
\endhead
\bottomrule\noalign{}
\endlastfoot
0. The audience is a group of clever physics Ph.D.~students who do not
know the solution; they must be able to learn from it and to teach it. &
Every term and symbol is defined in words before it is used; derivations
are given in steps that can be repeated by hand, and the text says how;
each result names the check that verifies it. & any section of a
manuscript or of ``dirac matrices deflating'' \\
1. All documentation, symbol and provenance files exist as .md, .tex and
.pdf, complete, precise and correct. & Each of the seven documents of
Section 4 exists in the three forms; the .tex and .pdf are built from
the .md (for the manuscripts: the .md from the .tex) by scripts, and the
build fails on an undefined reference or on a box wider than its space.
& the files; \texttt{code/scripts/build\_docs.sh},
\texttt{build\_manuscripts.py} \\
2. Push to the repository as soon as an object is updated. & Every
commit is carried at once to the author's folder and fast-forwarded
there. The push itself is done from the author's computer (Section 7,
item V3). & \texttt{git\ log} \\
3., 4. Working directory and repository. & The folder
\texttt{Revision\_Deflation\_CJT\_and\_Kohn-Sham} of the repository
\texttt{https://github.com/once-ere/spacetimes}, in the working
directory named by the rule. & \\
5. No symbol with two meanings; a MASTER provenance of the symbols,
promulgated and used. & Document \texttt{P1} and the table
\texttt{provenance/symbols\_master.csv}; the checks MS1 to MS7 (7 of 7
pass) verify the registry, the documents and the columns of all data
files against it. & \texttt{P1}, Section 5;
\texttt{data/symbol\_checks.txt} \\
6. Source code for every equation, result, figure and prediction; a
figure for each prediction; commented code, instructions and provenance
for every computed value; Rust preferred; the code executed. & One Rust
program (33 source files, 15007 lines, 68 unit tests, 18 commands)
computes everything; 30 figures are drawn from its data files; document
\texttt{P2} gives, for every table and figure, the command, the source
files and the data files, and prints the complete output of every check
(295 checks of the program, all passing). Everything was executed; the
reports in \texttt{data} are the outputs. Python is used for three
purposes only (Section 7, item V2). & \texttt{P2} \\
7. Never send the reader to another manuscript, document, provenance or
symbol file; every equation, symbol and definition in the file itself. &
See Section 2.2, first row. Program output printed in the documents uses
the expression ``the main project'' for the flat-spacetime part of
PROJECT spacetimes, whose program is used as a library and whose data
files several checks compare against; the documents that print such
output say so. & \texttt{data/self\_contained\_checks.txt} \\
8. Follow the instructions exactly, or explain in a provenance why not,
and pause. No shortcuts. & Sections 7 to 10 of this document. & \\
\end{longtable}

\hypertarget{what-the-folder-contains}{%
\section{4. What the folder contains}\label{what-the-folder-contains}}

\begin{longtable}[]{@{}
  >{\raggedright\arraybackslash}p{(\columnwidth - 2\tabcolsep) * \real{0.5000}}
  >{\raggedright\arraybackslash}p{(\columnwidth - 2\tabcolsep) * \real{0.5000}}@{}}
\toprule\noalign{}
\begin{minipage}[b]{\linewidth}\raggedright
document (each as .md, .tex, .pdf)
\end{minipage} & \begin{minipage}[b]{\linewidth}\raggedright
content
\end{minipage} \\
\midrule\noalign{}
\endhead
\bottomrule\noalign{}
\endlastfoot
\texttt{manuscript/16\ Component\ Dirac\ Fermion\ Condensate\ Singlet\ Superfluid\ on\ 4+4\ Deflating\ Spacetime}
& the manuscript for the 4+4 spacetime, in the style of Physical Review
D \\
\texttt{manuscript/16\ Component\ Dirac\ Fermion\ Condensate\ Singlet\ Superfluid\ on\ 7+1\ Deflating\ Spacetime}
& the manuscript for the 7+1 spacetime \\
\texttt{dirac\ matrices\ deflating} & geometry, the sixteen Dirac
matrices printed in full, spin connection, the Lagrangians with mass
term and self-interaction, the field equations, the algebra of the
self-interaction, the forms of the 4+4 model, solutions mode by mode;
complete reports of the checks \\
\texttt{provenance/P0\_task\_rules\_deviations\_decisions} & this
document \\
\texttt{provenance/P1\_master\_symbols\_spacetimes} & the MASTER
registry of symbols, the collisions and their resolution, the dictionary
to the author's 7+1 document, the machine checks \\
\texttt{provenance/P2\_calculations\_code\_and\_checks} & every
calculation: commands, instructions, data files, figures, complete
output of all checks, source files \\
\texttt{README} & the list of contents and the instructions for building
everything \\
\end{longtable}

\begin{longtable}[]{@{}
  >{\raggedright\arraybackslash}p{(\columnwidth - 2\tabcolsep) * \real{0.5000}}
  >{\raggedright\arraybackslash}p{(\columnwidth - 2\tabcolsep) * \real{0.5000}}@{}}
\toprule\noalign{}
\begin{minipage}[b]{\linewidth}\raggedright
other content
\end{minipage} & \begin{minipage}[b]{\linewidth}\raggedright
\end{minipage} \\
\midrule\noalign{}
\endhead
\bottomrule\noalign{}
\endlastfoot
\texttt{code/deflation-cjt-ks-rs} & the Rust program
\texttt{deflation\_cjt\_ks} \\
\texttt{code/scripts} & independent verification in Python (two scripts
and a test of them with deliberate errors); figures; the scripts that
build and check the documents \\
\texttt{data} & 111 data files written by the program, and the reports
of all checks \\
\texttt{figures} & 30 figures, each as .pdf and .png \\
\texttt{docs\_src} & the sources of all documents \\
\end{longtable}

The two manuscripts are built from one source
(\texttt{docs\_src/manuscript}): the text that the two share is written
once, and passages that differ are marked for one spacetime or the
other, so that the manuscripts cannot drift apart.

\hypertarget{what-was-done-in-order}{%
\section{5. What was done, in order}\label{what-was-done-in-order}}

\begin{longtable}[]{@{}
  >{\raggedright\arraybackslash}p{(\columnwidth - 2\tabcolsep) * \real{0.5000}}
  >{\raggedright\arraybackslash}p{(\columnwidth - 2\tabcolsep) * \real{0.5000}}@{}}
\toprule\noalign{}
\begin{minipage}[b]{\linewidth}\raggedright
commit
\end{minipage} & \begin{minipage}[b]{\linewidth}\raggedright
content
\end{minipage} \\
\midrule\noalign{}
\endhead
\bottomrule\noalign{}
\endlastfoot
\texttt{8511eeb} & the program with its 18 commands, all data files, the
figures, the two independent verifications and the test with deliberate
errors \\
\texttt{a232561} & the two manuscripts; \texttt{P1}; \texttt{P2} \\
\texttt{990d697} & correction of a failing symbol check (Section 12,
C-E12) \\
\texttt{8231b87} & correction of a failing number check and of an
unresolved placeholder; the script \texttt{check\_all.sh} (C-E13,
C-E14) \\
\texttt{56f9162} & the document ``dirac matrices deflating'';
corrections C-E15, C-E16 \\
the commit that contains this document & \texttt{P0} and README \\
\end{longtable}

\hypertarget{the-results-in-brief}{%
\section{6. The results in brief}\label{the-results-in-brief}}

Each statement is verified by the checks named in the manuscripts and in
\texttt{P2}. Numbers are in units of the fermion mass \(m\) unless said
otherwise.

\textbf{The Lagrangian and the field equation.} On the deflating
spacetimes the field \(\hat\psi=\sqrt{V_d}\,\psi\), weighted with the
proper volume \(V_d\) of the deflating coordinates, obeys for fields
that do not depend on \(x5,x6,x7\) the equation of flat spacetime in
which the mass is unchanged and \emph{every quartic coupling is divided
by \(V_d\)}: all couplings \(g_r\) grow as \(e^{\mathrm{gsum}\,x4}\),
with \(\mathrm{gsum}=\mathrm{g55}+\mathrm{g66}+\mathrm{g77}\).

\textbf{Nine couplings, four independent.} For anticommuting fields the
nine invariants obey \((1+\mathsf F)\,I=0\) with the Fierz matrix
\(\mathsf F\): only four combinations of the nine couplings are
physical. Uniform matter feels two of them, the Hartree--Fock couplings
\(G_S^{\rm HF}\) and \(G_V^{\rm HF}\); two particles at rest feel the
pair couplings, which take two values on the 28 pairs. A vector or
axial-vector coupling breaks the internal symmetry \(U(1)_E\) of fields
in ordinary space.

\textbf{Kohn--Sham.} The author's reference example on flat 7+1
spacetime is reproduced: saturation at the Fermi momentum 0.7566, mass
in matter 0.3481, binding energy per particle 0.0283. New: on the
deflating 7+1 spacetime matter that fills all momenta is self-bound only
in a window of the clock \(\tau_d=\mathrm{gsum}\,(x4-x4_{\rm ref})\),
between \(-0.675\) and 0.537 for equal rates, and the mass in matter
falls to half at the clock reading 0.983; in the sectors of fields that
do not depend on the deflating coordinates the thresholds of the
coupling for self-bound matter are computed in four and in three space
dimensions; in the covariant form of the 4+4 model matter polarizes
spontaneously between the two halves of the field.

\textbf{CJT.} For uniform matter the CJT effective potential, truncated
at the same order, has the Kohn--Sham state as its stationary point: the
two methods agree to \(1.6\times10^{-12}\) on 30 states, and the CJT
potential shows in addition which of them are global minima. With the
filled states of negative energy and a momentum cutoff \(\Lambda\) the
vacuum acquires a mass above a critical coupling; the author's coupling
exceeds it for cutoffs above 0.6665. A unified functional contains the
contact couplings, a bare mass and the Vorton-Scalar; a bare mass gives
the phase mode of the Vorton-Scalar a mass.

\textbf{Pairing.} The exchange of the radial mode of the Vorton-Scalar
is a scalar contact coupling with a finite range; with contact couplings
added, the gap of the hidden-axis pair changes by large factors (a
factor 12.0 for the scalar attraction of the example).

\textbf{Gravity.} With the matter as the source of Einstein--Lovelock
equations for two scale factors, the author's reference evolution is
reproduced to \(4.8\times10^{-8}\). The exponential deflation with equal
rates is a solution of Einstein--Gauss--Bonnet gravity without matter,
and it is unstable there: matter makes the deflating coordinates deflate
faster than exponentially, and in four of eight evolutions the rates
grow without bound.

\textbf{The earlier results}, on the modes, the clock, the pairing
through the Vorton-Scalar, the bound pair, the quench, the effective
action along the clock, the cosmology and the Higgs hypothesis, were
computed again by the commands of the new program and are stated in the
manuscripts with the changes that the mass term and the self-interaction
bring (manuscripts, Sec. XVIII B).

\hypertarget{deviations-from-the-instructions-and-the-reasons}{%
\section{7. Deviations from the instructions, and the
reasons}\label{deviations-from-the-instructions-and-the-reasons}}

\textbf{V1. The five supplements of the first revision are not repeated
as separate documents.} Hard Rule 7 forbids a document to send its
reader to another, and a supplement is a document to which a manuscript
sends its reader. The content of the supplements is therefore inside
each manuscript:

\begin{longtable}[]{@{}
  >{\raggedright\arraybackslash}p{(\columnwidth - 2\tabcolsep) * \real{0.5000}}
  >{\raggedright\arraybackslash}p{(\columnwidth - 2\tabcolsep) * \real{0.5000}}@{}}
\toprule\noalign{}
\begin{minipage}[b]{\linewidth}\raggedright
supplement of the first revision
\end{minipage} & \begin{minipage}[b]{\linewidth}\raggedright
where its content is in each manuscript of this revision
\end{minipage} \\
\midrule\noalign{}
\endhead
\bottomrule\noalign{}
\endlastfoot
Vorton-Scalar and superfluid & Secs. VI, VIII, X, XI, XIV \\
effective action along the clock & Sec. XII \\
cosmology & Sec. XVI A to D, F, I to K \\
microwave background & Sec. XVI E, G, H \\
Higgs hypothesis & Sec. XVII \\
\end{longtable}

The derivations are in the sections named. Of the tables of the
supplements, the manuscripts print those that carry the argument and
quote numbers from the others; the complete tables that are only quoted
are generated in \texttt{provenance/generated/tables\_cjt\_ks.md} and
their data files are listed in \texttt{P2}, Section 6. Printing them as
an appendix of each manuscript is the last step of Section 10. If the
author wants the supplements as separate, self-contained documents as
well, that is a matter of assembling existing material; it has not been
done because it would repeat, five times over, the definitions that each
manuscript already contains.

\textbf{V2. Python is used beside Rust}, for three purposes: (i) the
independent verification, which must not share code with the program
that it verifies and therefore uses another language and other libraries
(computer algebra, a general-purpose integrator, other special
functions); (ii) the figures; (iii) the scripts that assemble, typeset
and check the documents. No result of a manuscript comes from Python
alone. Mathematica was not available; the numbers of the author's
notebook that his 7+1 document prints are reproduced by the Rust program
and by the independent Python script.

\textbf{V3. The push to GitHub is done from the author's computer.} The
work is carried out in a workspace without the author's credentials.
Each commit is transferred at once to his folder as a git bundle and
fast-forwarded there (the folder is then one commit ahead of the remote
until he pushes).

\textbf{V4. The test with deliberate errors does not cover the whole of
the first independent script.} It introduces 36 errors (20 in the script
for the self-interaction, the Kohn-Sham and CJT calculations, pairing
and gravity, 16 in the cosmological part of the other script); all are
detected. The part of the first script that verifies the geometry, the
Dirac operators and the modes has no such test.

\textbf{V5. The titles of the manuscripts} are those of the first
revision; the task names no titles.

\textbf{V6. An external Einstein--Boltzmann code was not used as a
cross-check.} The spectra of the microwave background are computed with
the program of this project, whose version for static extra coordinates
had been validated in the flat-spacetime part of the project; no
external code was built or run in this sub-project.

\hypertarget{decisions-that-are-the-authors}{%
\section{8. Decisions that are the
author's}\label{decisions-that-are-the-authors}}

\textbf{D1. The form of the self-interaction on the 4+4 spacetime.} This
is the point at which the author's 7+1 construction does not carry over
mechanically. The covariant 4+4 Lagrangian
\(\tilde\psi(\gamma^\mu\nabla_\mu-m)\psi\), with
\(\tilde\psi=\psi^\dagger C^-\), is indefinite: for fields that do not
depend on \(x5,x6,x7\) it is the ordinary Dirac Lagrangian on one half
of the field and minus it on the other half (the halves are the
eigenspaces of \(W_{1644}=iC^-\gamma^4\)). For a free field that is
harmless. With a self-interaction the two halves are no longer
equivalent, and two definitions are possible:

\begin{itemize}
\tightlist
\item
  \emph{Q, the quantized-sector form (adopted):} bilinears with
  \(\bar\psi=\psi^\dagger\beta\), \(\beta=-i\gamma^4\). All results of
  the 7+1 model then hold on 4+4 without change.
\item
  \emph{C, the covariant form with the positive kinetic term (reported
  beside Q):} the author's invariants with \(\tilde\psi\). It equals Q
  with \(W_{1644}\) inserted in every bilinear; particles in different
  halves interact with the opposite sign; the hidden-axis pair consists
  of particles in different halves, so that a scalar attraction repels
  in that channel; the condition for chiral symmetry changes.
\end{itemize}

The author was notified of this choice when it was made and again when
the documents were complete. If he prefers C, the text can adopt it:
both are computed.

\textbf{D2. The bare mass.} With the Vorton-Scalar, a bare mass \(m_D\)
breaks the chiral symmetry explicitly and gives the phase mode a mass.
Results that rest on a massless phase mode (the stepped fluid of the
cosmological part) presuppose \(m_D=0\) or a mass of the phase mode
below the temperatures concerned. Which is meant is the author's
decision.

\textbf{D3. The couplings of the reference example.} See Section 9, F2:
are \(-10^5\) and \(+10^4\) the couplings \(g_0\), \(g_1\) of the
Lagrangian or the couplings \(G_S\), \(G_V\) of the mean-field energy?

\textbf{D4. The symbols of the metric.} The metric of this sub-project
is written \(g^{(4+4)}\), \(g^{(7+1)}\) in the author's sign convention
(space-like components positive), and the rates keep his names
\(\mathrm{g55},\mathrm{g66},\mathrm{g77}\), upright. The names of the
first revision for a metric of the other sign convention are retired.

\textbf{D5. Open from the first revision:} the assumption about the
perturbations of ordinary space and its closures; the nature of whatever
sustains the deflation; the choice between Einstein gravity with such an
agent and Einstein--Gauss--Bonnet gravity.

\hypertarget{findings-that-concern-the-authors-own-documents}{%
\section{9. Findings that concern the author's own
documents}\label{findings-that-concern-the-authors-own-documents}}

\textbf{F1. Only four of the nine couplings are independent.} The
identity \((1+\mathsf F)I=0\) holds for anticommuting fields with the
Fierz matrix of his document. A tensor coupling \(g_2,\dots,g_6\) is not
a new parameter.

\textbf{F2. The couplings of the Kohn--Sham example are not the
Hartree--Fock couplings of \(g_0=-10^5\), \(g_1=+10^4\).} His document
states the Hartree--Fock couplings of these values correctly as
\(-98750\) and \(+20000\) (check \texttt{self-X1}), and its Kohn--Sham
calculation uses \(G_S=-10^5\), \(G_V=+10^4\). With the first pair,
matter on flat 7+1 spacetime is not self-bound: its energy per particle
has only a local minimum above the mass. The numbers of his document are
reproduced with the second pair. If the mean-field couplings are meant
to be the Hartree--Fock couplings, the Lagrangian couplings that give
his example are \(g_0=-105341.6\) and \(g_1=+2484.5\) in the same units.

\textbf{F3. The coupling of the example exceeds the critical coupling of
the vacuum} when the filled states of negative energy are included with
a cutoff above 0.6665: the Kohn--Sham calculation without those states
presupposes a cutoff below that value.

\textbf{F4. The deflating spacetime is an unstable solution} of the
gravity in which it is a vacuum solution, and in Einstein gravity it
needs a source without energy density and with a tension; Kohn--Sham
matter cannot be that source.

\textbf{F5. A vector or axial-vector coupling breaks \(U(1)_E\)}, the
second phase symmetry of the fields in ordinary space, on which the
hypothesis about the Higgs field rests.

\hypertarget{what-has-not-been-done}{%
\section{10. What has not been done}\label{what-has-not-been-done}}

\begin{enumerate}
\def\labelenumi{\arabic{enumi}.}
\tightlist
\item
  \emph{Complete tables as an appendix of the manuscripts.} Twenty-six
  tables whose numbers the manuscripts quote in part (thermodynamic
  ratios and Ginzburg--Landau coefficients of the paired state; melting
  temperatures; the vacuum potential; the drift of quantities; the loss
  of rest mass of bound pairs; details of the microwave-background
  calculation and of the fits to mock data; the limits on superfluid
  dark matter; the internal symmetry item by item) are generated but not
  yet printed in the manuscripts.
\item
  \emph{Matter beyond the mean field; pairing inside dense self-bound
  matter; a self-consistent solution with the created particles acting
  on the metric.} Stated as not established in the manuscripts, Sec.
  XVIII D.
\item
  \emph{The covariant form C of the 4+4 model} is computed for the
  couplings, for uniform matter and for pairs at rest, not for every
  later result (gap along the clock, gravity).
\item
  \emph{A test with deliberate errors for the geometric part of the
  first independent script} (Section 7, item V4).
\item
  \emph{The paused stages of the flat-spacetime part of the project}
  were not touched, as instructed.
\end{enumerate}

\hypertarget{inputs}{%
\section{11. Inputs}\label{inputs}}

\begin{longtable}[]{@{}
  >{\raggedright\arraybackslash}p{(\columnwidth - 2\tabcolsep) * \real{0.5000}}
  >{\raggedright\arraybackslash}p{(\columnwidth - 2\tabcolsep) * \real{0.5000}}@{}}
\toprule\noalign{}
\begin{minipage}[b]{\linewidth}\raggedright
input
\end{minipage} & \begin{minipage}[b]{\linewidth}\raggedright
what was taken from it
\end{minipage} \\
\midrule\noalign{}
\endhead
\bottomrule\noalign{}
\endlastfoot
the author's task specification for this sub-project (private; not in
the repository) & the hypotheses and the rules of Sections 2 and 3 \\
his document \texttt{KOHN\_SHAM\_GRAVITY} (folder \texttt{7+1} on his
computer; its .tex source was read in full) & the definition of
\(\mathcal L_{\rm self}\), of \(\Theta^K\), \(I_r\), \(g_r\); the Fierz
matrix; the Kohn--Sham equations and the reference numbers; the
Einstein--Lovelock equations for two scale factors and the reference
evolution \\
his document \texttt{DIRAC16COMPLEX} (same folder) & the critical
coupling of the vacuum in seven space dimensions \\
the notebooks and papers of the project, as recorded in the
flat-spacetime part & the two sets of real \(16\times16\) Dirac matrices
(files \texttt{data/gamma1644.txt}, \texttt{data/gamma1671.txt} at the
top of the repository); the conventions of the fields \\
the folder \texttt{Revision\_Deflation} (the first revision) & the
calculations that are repeated here; it was not changed \\
\end{longtable}

The author's files were read, not changed, and none of them is copied
into the repository.

\hypertarget{log-of-errors}{%
\section{12. Log of errors}\label{log-of-errors}}

Every error made during this sub-project is recorded here, including
errors that were corrected before anything was committed. For the errors
C-E1 to C-E11, which belong to that kind, the record kept during the
work names the error and what exposed it; an entry says no more than the
record contains.

\begin{longtable}[]{@{}
  >{\raggedright\arraybackslash}p{(\columnwidth - 6\tabcolsep) * \real{0.2500}}
  >{\raggedright\arraybackslash}p{(\columnwidth - 6\tabcolsep) * \real{0.2500}}
  >{\raggedright\arraybackslash}p{(\columnwidth - 6\tabcolsep) * \real{0.2500}}
  >{\raggedright\arraybackslash}p{(\columnwidth - 6\tabcolsep) * \real{0.2500}}@{}}
\toprule\noalign{}
\begin{minipage}[b]{\linewidth}\raggedright
No.
\end{minipage} & \begin{minipage}[b]{\linewidth}\raggedright
what was wrong
\end{minipage} & \begin{minipage}[b]{\linewidth}\raggedright
how it was found
\end{minipage} & \begin{minipage}[b]{\linewidth}\raggedright
correction
\end{minipage} \\
\midrule\noalign{}
\endhead
\bottomrule\noalign{}
\endlastfoot
C-E1 & The first formula for the pair couplings acted on pair states
without projecting on the states of positive energy and gave other
integers. & Comparison with the exact diagonalization for two particles
at rest. & The projected map; checks \texttt{71s-S9},
\texttt{44q-S9}. \\
C-E2 & An early text called the author's values \(G_S\), \(G_V\) of the
Kohn--Sham example the Hartree--Fock couplings of \(g_0\), \(g_1\) with
the same numbers. & The exact Hartree--Fock formula (check
\texttt{self-X1}). & The texts distinguish the two; finding F2. \\
C-E3 & In the first version of the check of the mass of the phase mode
(\texttt{uni-U3}) a deviation appeared that was an artefact of the
momentum cutoff. & The check failed. & Both sides of the comparison are
evaluated with the same cutoff. \\
C-E4 & In the history of deflating matter (\texttt{ks-K5}) a
minimization ran without bounds. & Unphysical values in the output. &
The search is bounded. \\
C-E5 & An angular average of the pairing amplitudes evaluated a function
outside its domain. & The program stopped. & Corrected. \\
C-E6 & The tolerance of an independent comparison of two integrations of
the deflating spacetime with matter (\texttt{Icks-G5}) was set tighter
than the accuracy of the comparison. & The check failed although both
integrations were correct. & The tolerance follows from the measured
accuracy, which the check states. \\
C-E7 & Labels of several new figures collided with curves, and some used
the names of the program or of the author's document instead of the
registered symbols. & Looking at each figure. & Figures redrawn. \\
C-E8 & The registry of symbols at first contained collisions introduced
by this sub-project: one letter for a velocity dispersion and for a
self-energy, one for a pair state and for a pair correlation, and two
definitions of \(x\). & The checks of the registry and reading the
manuscripts. & New symbols \(\sigma_v\), \(\Phi(\vec p)\); one
definition of \(x\). \\
C-E9 & A generated table contained section numbers typed by hand. &
Reading the built manuscript. & The table carries references that the
typesetting resolves. \\
C-E10 & The description of the test with deliberate errors repeated a
statement that belonged to a test of the first revision. & Comparison
with the report of the script of this folder. & The text states what the
report of this folder says. \\
C-E11 & A clause about the degeneracy of pair states in the section on
the Higgs hypothesis had not been verified. & Review of the section
against the checks. & Removed. \\
C-E12 & \textbf{A failing check was committed.} The message of commit
\texttt{a232561} says that the symbol checks pass \texttt{7\ of\ 7}; the
report in that commit says \texttt{6\ of\ 7}. With the documents
\texttt{P1} and \texttt{P2} in the folder, check MS4 failed: two
generated texts wrote the rates as plain text. & Reading the whole
report instead of its last lines, one step later. & Commit
\texttt{990d697}: both texts write the rates upright. \\
C-E13 & \textbf{A second failing check was committed.} Since \texttt{P2}
was added, the check of the numbers failed (it did not recognize
\texttt{Section\ 5.11} as a section number, and two run times had no
entry in its lists), and \texttt{P2} printed the failing report. & As
C-E12. & Commit \texttt{8231b87}. \\
C-E14 & \texttt{P2} printed a caption with an unresolved placeholder. &
The new script \texttt{check\_all.sh}. & Commit \texttt{8231b87}. \\
C-E15 & The running header of \texttt{P1} and \texttt{P2} named the
folder of the first revision. & Looking at the typeset pages. & Commit
\texttt{56f9162}. \\
C-E16 & In the figure of the creation of particles the labels of two
axes touched, and two legends lay on curves; the figure is in the 7+1
manuscript. & Looking at the typeset page of the new document. & Commit
\texttt{56f9162}: labels on two lines, legends moved, manuscript
rebuilt. \\
C-E17 & A sentence drafted for ``dirac matrices deflating'' took the
roots of the constraint of the gravitational equations for rates of
growth. & Reading the statement of check \texttt{gr-G5} before the
document was committed. & The document says what the check says. \\
\end{longtable}

\textbf{The cause of C-E12 to C-E14, and the remedy.} After long builds
only the last lines of the outputs were read, and a verdict further up
was missed. The script \texttt{code/scripts/check\_all.sh} now runs the
checks of the documents, reads the verdict of every report in
\texttt{data}, verifies that no built document contains an unresolved
placeholder and that every report printed in a document is the current
one, and prints \texttt{ALL\ PASS} only then. Nothing is committed
without it.

\hypertarget{the-symbols-used-in-this-document}{%
\section{13. The symbols used in this
document}\label{the-symbols-used-in-this-document}}

\begin{longtable}[]{@{}
  >{\raggedright\arraybackslash}p{(\columnwidth - 2\tabcolsep) * \real{0.5000}}
  >{\raggedright\arraybackslash}p{(\columnwidth - 2\tabcolsep) * \real{0.5000}}@{}}
\toprule\noalign{}
\begin{minipage}[b]{\linewidth}\raggedright
symbol
\end{minipage} & \begin{minipage}[b]{\linewidth}\raggedright
meaning
\end{minipage} \\
\midrule\noalign{}
\endhead
\bottomrule\noalign{}
\endlastfoot
\(x0,\dots,x7\); \(xt\) & the eight coordinates (\(x4\): ordinary time);
a deflating coordinate, \(t=5,6,7\) \\
\(\mathrm{g55}\), \(\mathrm{g66}\), \(\mathrm{g77}\);
\(\mathrm{rate}_t\); \(\mathrm{gsum}\) & the three rates of deflation
(the author's names; upright, never with a subscript); the rate of
\(xt\); their sum \\
\(\mathrm{const}_t\) & the positive constants of the hypothesis, one for
each deflating coordinate \\
\(g^{(4+4)}_{\mu\nu}\), \(g^{(7+1)}_{\mu\nu}\) & the metrics of the two
deflating spacetimes, space-like components positive \\
\(V_d\) & \(\sqrt{\lvert\det g\rvert}\), the proper volume of a unit
coordinate cell of \(x5,x6,x7\); proportional to
\(e^{-\mathrm{gsum}\,x4}\) \\
\(\tau_d\) & the deflation clock, \(\mathrm{gsum}\,(x4-x4_{\rm ref})\),
with a reference time \(x4_{\rm ref}\) \\
\(\psi\); \(\hat\psi\) & the sixteen-component Dirac field; the weighted
field \(\sqrt{V_d}\,\psi\) \\
\(\gamma^A\); \(\gamma^\mu\); \(\nabla_\mu\) & the eight real
\(16\times16\) Dirac matrices of an orthonormal frame, \(A=0,\dots,7\);
the same with a coordinate index; the covariant derivative of a
spinor \\
\(\beta\); \(\bar\psi\) & \(\gamma^4\) on 7+1 and \(-i\gamma^4\) on 4+4;
the adjoint \(\psi^\dagger\beta\) \\
\(C^-\); \(\tilde\psi\); \(W_{1644}\) & on 4+4: the matrix
\(\gamma^0\gamma^1\gamma^2\gamma^3\); the covariant adjoint
\(\psi^\dagger C^-\); the matrix \(iC^-\gamma^4\), with eigenvalues
\(\pm1\) \\
\(m\); \(m_D\) & the mass of the fermion; the part of it that is a bare
mass in the model with the Vorton-Scalar \\
\(\mathcal L_{\rm kin}\), \(\mathcal L_{\rm mass}\),
\(\mathcal L_{\rm self}\) & the three parts of the Lagrangian: kinetic
term, mass term, self-interaction \\
\(K\); \(r\); \(\Theta^K\), \(\Theta_K\) & a strictly increasing list of
indices out of \(0,\dots,7\); its length, the rank; the product of the
Dirac matrices of the list, with a phase that makes the bilinears real,
and its inverse \\
\(I_r\); \(g_r\) & the invariant of rank \(r\); the nine couplings,
\(r=0,\dots,8\) (\(g_0\) scalar, \(g_1\) vector) \\
\(\mathsf F\) & the Fierz matrix, a \(9\times9\) matrix of rational
numbers with \(\mathsf F^2=1\) \\
\(G_S\), \(G_V\); \(G_S^{\rm HF}\), \(G_V^{\rm HF}\) & the scalar and
the vector coupling of the mean-field energy of uniform matter; their
values from the nine couplings in the Hartree--Fock approximation \\
\(k_F\) & the Fermi momentum \\
\(\Lambda\) & a momentum cutoff \\
\(U(1)_E\) & a phase symmetry of the fields that depend on
\(x1,x2,x3,x4\) only \\
Vorton-Scalar & the complex scalar field of the model, coupled to
\(\bar\psi\psi\); its radial mode and its phase mode are its two kinds
of excitation \\
\end{longtable}

\end{document}
